\section{Proj of a graded ring}
\label{section-proj}

\noindent
Let $S$ be a graded ring.
A {\it homogeneous ideal} is simply an ideal
$I \subset S$ which is also a graded submodule of $S$.
Equivalently, it is an ideal generated by homogeneous elements.
Equivalently, if $f \in I$ and
$$
f = f_0 + f_1 + \ldots + f_n
$$
is the decomposition of $f$ into homogeneous parts in $S$ then $f_i \in I$
for each $i$. To check that a homogeneous ideal $\mathfrak p$
is prime it suffices to check that if $ab \in \mathfrak p$
with $a, b$ homogeneous then either $a \in \mathfrak p$ or
$b \in \mathfrak p$.

\begin{definition}
\label{definition-proj}
Let $S$ be a graded ring.
We define $\text{Proj}(S)$ to be the set of homogeneous
prime ideals $\mathfrak p$ of $S$ such that
$S_{+} \not \subset \mathfrak p$.
The set $\text{Proj}(S)$ is a subset of $\Spec(S)$
and we endow it with the induced topology.
The topological space $\text{Proj}(S)$ is called the
{\it homogeneous spectrum} of the graded ring $S$.
\end{definition}

\noindent
Note that by construction there is a continuous map
$$
\text{Proj}(S) \longrightarrow \Spec(S_0).
$$

\medskip\noindent
Let $S = \oplus_{d \geq 0} S_d$ be a graded ring.
Let $f\in S_d$ and assume that $d \geq 1$.
We define $S_{(f)}$ to be the subring of $S_f$
consisting of elements of the form $r/f^n$ with $r$ homogeneous and
$\deg(r) = nd$. If $M$ is a graded $S$-module,
then we define the $S_{(f)}$-module $M_{(f)}$ as the
sub module of $M_f$ consisting of elements of
the form $x/f^n$ with $x$ homogeneous of degree $nd$.

\begin{lemma}
\label{lemma-Z-graded}
Let $S$ be a $\mathbf{Z}$-graded ring containing a homogeneous
invertible element of positive degree. Then the set
$G \subset \Spec(S)$ of $\mathbf{Z}$-graded primes of $S$
(with induced topology) maps homeomorphically to $\Spec(S_0)$.
\end{lemma}

\begin{proof}
First we show that the map is a bijection by constructing an inverse.
Let $f \in S_d$, $d > 0$ be invertible in $S$.
If $\mathfrak p_0$ is a prime of $S_0$, then $\mathfrak p_0S$
is a $\mathbf{Z}$-graded ideal of $S$ such that
$\mathfrak p_0S \cap S_0 = \mathfrak p_0$. And if $ab \in \mathfrak p_0S$
with $a$, $b$ homogeneous, then
$a^db^d/f^{\deg(a) + \deg(b)} \in \mathfrak p_0$.
Thus either $a^d/f^{\deg(a)} \in \mathfrak p_0$ or
$b^d/f^{\deg(b)} \in \mathfrak p_0$, in other words either
$a^d \in \mathfrak p_0S$ or $b^d \in \mathfrak p_0S$.
It follows that $\sqrt{\mathfrak p_0S}$ is a $\mathbf{Z}$-graded
prime ideal of $S$ whose intersection with $S_0$ is $\mathfrak p_0$.

\medskip\noindent
To show that the map is a homeomorphism we show that
the image of $G \cap D(g)$ is open. If $g = \sum g_i$
with $g_i \in S_i$, then by the above $G \cap D(g)$
maps onto the set $\bigcup D(g_i^d/f^i)$ which is open.
\end{proof}

\noindent
For $f \in S$ homogeneous of degree $> 0$ we define
$$
D_{+}(f) = \{ \mathfrak p \in \text{Proj}(S) \mid f \not\in \mathfrak p \}.
$$
Finally, for a homogeneous ideal $I \subset S$ we define
$$
V_{+}(I) = \{ \mathfrak p \in \text{Proj}(S) \mid I \subset \mathfrak p \}.
$$
We will use more generally the notation $V_{+}(E)$ for any
set $E$ of homogeneous elements $E \subset S$.

\begin{lemma}[Topology on Proj]
\label{lemma-topology-proj}
Let $S = \oplus_{d \geq 0} S_d$ be a graded ring.
\begin{enumerate}
\item The sets $D_{+}(f)$ are open in $\text{Proj}(S)$.
\item We have $D_{+}(ff') = D_{+}(f) \cap D_{+}(f')$.
\item Let $g = g_0 + \ldots + g_m$ be an element
of $S$ with $g_i \in S_i$. Then
$$
D(g) \cap \text{Proj}(S) =
(D(g_0) \cap \text{Proj}(S))
\cup
\bigcup\nolimits_{i \geq 1} D_{+}(g_i).
$$
\item
Let $g_0\in S_0$ be a homogeneous element of degree $0$. Then
$$
D(g_0) \cap \text{Proj}(S)
=
\bigcup\nolimits_{f \in S_d, \ d\geq 1} D_{+}(g_0 f).
$$
\item The open sets $D_{+}(f)$ form a
basis for the topology of $\text{Proj}(S)$.
\item Let $f \in S$ be homogeneous of positive degree.
The ring $S_f$ has a natural $\mathbf{Z}$-grading.
The ring maps $S \to S_f \leftarrow S_{(f)}$ induce
homeomorphisms
$$
D_{+}(f)
\leftarrow
\{\mathbf{Z}\text{-graded primes of }S_f\}
\to
\Spec(S_{(f)}).
$$
\item There exists an $S$ such that $\text{Proj}(S)$ is not
quasi-compact.
\item The sets $V_{+}(I)$ are closed.
\item Any closed subset $T \subset \text{Proj}(S)$ is of
the form $V_{+}(I)$ for some homogeneous ideal $I \subset S$.
\item For any graded ideal $I \subset S$ we have
$V_{+}(I) = \emptyset$ if and only if $S_{+} \subset \sqrt{I}$.
\end{enumerate}
\end{lemma}

\begin{proof}
Since $D_{+}(f) = \text{Proj}(S) \cap D(f)$, these sets are open.
This proves (1). Also (2) follows as $D(ff') = D(f) \cap D(f')$.
Similarly the sets $V_{+}(I) = \text{Proj}(S) \cap V(I)$
are closed. This proves (8).

\medskip\noindent
Suppose that $T \subset \text{Proj}(S)$ is closed.
Then we can write $T = \text{Proj}(S) \cap V(J)$ for some
ideal $J \subset S$. By definition of a homogeneous ideal
if $g \in J$, $g = g_0 + \ldots + g_m$
with $g_d \in S_d$ then $g_d \in \mathfrak p$ for all
$\mathfrak p \in T$. Thus, letting $I \subset S$
be the ideal generated by the homogeneous parts of the elements
of $J$ we have $T = V_{+}(I)$. This proves (9).

\medskip\noindent
The formula for $\text{Proj}(S) \cap D(g)$, with $g \in S$ is direct
from the definitions. This proves (3).
Consider the formula for $\text{Proj}(S) \cap D(g_0)$.
The inclusion of the right hand side in the left hand side is
obvious. For the other inclusion, suppose $g_0 \not \in \mathfrak p$
with $\mathfrak p \in \text{Proj}(S)$. If all $g_0f \in \mathfrak p$
for all homogeneous $f$ of positive degree, then we see that
$S_{+} \subset \mathfrak p$ which is a contradiction. This gives
the other inclusion. This proves (4).

\medskip\noindent
The collection of opens $D(g) \cap \text{Proj}(S)$
forms a basis for the topology since the standard opens
$D(g) \subset \Spec(S)$ form a basis for the topology on
$\Spec(S)$. By the formulas above we can express
$D(g) \cap \text{Proj}(S)$ as a union of opens $D_{+}(f)$.
Hence the collection of opens $D_{+}(f)$ forms a basis for the topology
also. This proves (5).

\medskip\noindent
Proof of (6). First we note that $D_{+}(f)$ may be identified
with a subset (with induced topology) of $D(f) = \Spec(S_f)$
via Lemma \ref{lemma-standard-open}. Note that the ring
$S_f$ has a $\mathbf{Z}$-grading. The homogeneous elements are
of the form $r/f^n$ with $r \in S$ homogeneous and have
degree $\deg(r/f^n) = \deg(r) - n\deg(f)$. The subset
$D_{+}(f)$ corresponds exactly to those prime ideals
$\mathfrak p \subset S_f$ which are $\mathbf{Z}$-graded ideals
(i.e., generated by homogeneous elements). Hence we have to show that
the set of $\mathbf{Z}$-graded prime ideals of $S_f$ maps homeomorphically
to $\Spec(S_{(f)})$. This follows from Lemma \ref{lemma-Z-graded}.

\medskip\noindent
Let $S = \mathbf{Z}[X_1, X_2, X_3, \ldots]$ with grading such that
each $X_i$ has degree $1$. Then it is easy to see that
$$
\text{Proj}(S) = \bigcup\nolimits_{i = 1}^\infty D_{+}(X_i)
$$
does not have a finite refinement. This proves (7).

\medskip\noindent
Let $I \subset S$ be a graded ideal.
If $\sqrt{I} \supset S_{+}$ then $V_{+}(I) = \emptyset$ since
every prime $\mathfrak p \in \text{Proj}(S)$ does not contain
$S_{+}$ by definition. Conversely, suppose that
$S_{+} \not \subset \sqrt{I}$. Then we can find an element
$f \in S_{+}$ such that $f$ is not nilpotent modulo $I$.
Clearly this means that one of the homogeneous parts of $f$
is not nilpotent modulo $I$, in other words we may (and do)
assume that $f$ is homogeneous. This implies that
$I S_f \not = S_f$, in other words that $(S/I)_f$ is not
zero. Hence $(S/I)_{(f)} \not = 0$ since it is a ring
which maps into $(S/I)_f$. Pick a prime
$\mathfrak q \subset (S/I)_{(f)}$. This corresponds to
a graded prime of $S/I$, not containing the irrelevant ideal
$(S/I)_{+}$. And this in turn corresponds to a graded prime
ideal $\mathfrak p$ of $S$, containing $I$ but not containing $S_{+}$
as desired. This proves (10) and finishes the proof.
\end{proof}

\begin{example}
\label{example-proj-polynomial-ring-1-variable}
Let $R$ be a ring. If $S = R[X]$ with $\deg(X) = 1$, then the natural map
$\text{Proj}(S) \to \Spec(R)$ is a bijection and in fact a homeomorphism.
Namely, suppose $\mathfrak p \in \text{Proj}(S)$. Since
$S_{+} \not \subset \mathfrak p$ we see that $X \not \in \mathfrak p$.
Thus if $aX^n \in \mathfrak p$ with $a \in R$ and $n > 0$, then
$a \in \mathfrak p$. It follows that $\mathfrak p = \mathfrak p_0S$
with $\mathfrak p_0 = \mathfrak p \cap R$.
\end{example}

\noindent
If $\mathfrak p \in \text{Proj}(S)$, then we
define $S_{(\mathfrak p)}$ to be the ring whose
elements are fractions $r/f$ where $r, f \in S$ are homogeneous
elements of the same degree such that $f \not\in \mathfrak p$.
As usual we say $r/f = r'/f'$ if and only if there exists
some $f'' \in S$ homogeneous, $f'' \not \in \mathfrak p$ such
that $f''(rf' - r'f) = 0$.
Given a graded $S$-module $M$ we let
$M_{(\mathfrak p)}$ be the $S_{(\mathfrak p)}$-module
whose elements are fractions $x/f$ with $x \in M$
and $f \in S$ homogeneous of the same degree such that
$f \not \in \mathfrak p$. We say $x/f = x'/f'$
if and only if there exists some $f'' \in S$ homogeneous,
$f'' \not \in \mathfrak p$ such that $f''(xf' - x'f) = 0$.

\begin{lemma}
\label{lemma-proj-prime}
Let $S$ be a graded ring. Let $M$ be a graded $S$-module.
Let $\mathfrak p$ be an element of $\text{Proj}(S)$.
Let $f \in S$ be a homogeneous element of positive degree
such that $f \not \in \mathfrak p$, i.e., $\mathfrak p \in D_{+}(f)$.
Let $\mathfrak p' \subset S_{(f)}$ be the element of
$\Spec(S_{(f)})$ corresponding to $\mathfrak p$ as in
Lemma \ref{lemma-topology-proj}. Then
$S_{(\mathfrak p)} = (S_{(f)})_{\mathfrak p'}$
and compatibly
$M_{(\mathfrak p)} = (M_{(f)})_{\mathfrak p'}$.
\end{lemma}

\begin{proof}
We define a map $\psi : M_{(\mathfrak p)} \to (M_{(f)})_{\mathfrak p'}$.
Let $x/g \in M_{(\mathfrak p)}$. We set
$$
\psi(x/g) = (x g^{\deg(f) - 1}/f^{\deg(x)})/(g^{\deg(f)}/f^{\deg(g)}).
$$
This makes sense since $\deg(x) = \deg(g)$ and since
$g^{\deg(f)}/f^{\deg(g)} \not \in \mathfrak p'$.
We omit the verification that $\psi$ is well defined, a module map
and an isomorphism. Hint: the inverse sends $(x/f^n)/(g/f^m)$ to
$(xf^m)/(g f^n)$.
\end{proof}

\noindent
Here is a graded variant of Lemma \ref{lemma-silly}.

\begin{lemma}
\label{lemma-graded-silly}
Suppose $S$ is a graded ring, $\mathfrak p_i$, $i = 1, \ldots, r$
homogeneous prime ideals and $I \subset S_{+}$ a graded ideal.
Assume $I \not\subset \mathfrak p_i$ for all $i$. Then there
exists a homogeneous element $x\in I$ of positive degree such
that $x\not\in \mathfrak p_i$ for all $i$.
\end{lemma}

\begin{proof}
We may assume there are no inclusions among the $\mathfrak p_i$.
The result is true for $r = 1$. Suppose the result holds for $r - 1$.
Pick $x \in I$ homogeneous of positive degree such that
$x \not \in \mathfrak p_i$ for all $i = 1, \ldots, r - 1$.
If $x \not\in \mathfrak p_r$ we are done. So assume $x \in \mathfrak p_r$.
If $I \mathfrak p_1 \ldots \mathfrak p_{r-1} \subset \mathfrak p_r$
then $I \subset \mathfrak p_r$ a contradiction.
Pick $y \in I\mathfrak p_1 \ldots \mathfrak p_{r-1}$ homogeneous
and $y \not \in \mathfrak p_r$. Then $x^{\deg(y)} + y^{\deg(x)}$ works.
\end{proof}

\begin{lemma}
\label{lemma-smear-out}
Let $S$ be a graded ring.
Let $\mathfrak p \subset S$ be a prime.
Let $\mathfrak q$ be the homogeneous ideal of $S$ generated by the
homogeneous elements of $\mathfrak p$. Then $\mathfrak q$ is a
prime ideal of $S$.
\end{lemma}

\begin{proof}
To prove that $\mathfrak q$ is prime, it suffices to check that
if $f, g \in S$ are homogeneous and $fg \in \mathfrak q$, then
either $f$ or $g$ in $\mathfrak q$. Then $fg \in \mathfrak p$
because $\mathfrak p \subset \mathfrak q$. Since $\mathfrak p$
is prime we see that either $f \in \mathfrak p$ or $g \in \mathfrak p$.
Since $f$ and $g$ are homogeneous, it then is clear that either
$f \in \mathfrak q$ or $g \in \mathfrak q$.
\end{proof}

\begin{lemma}
\label{lemma-graded-ring-minimal-prime}
Let $S$ be a graded ring.
\begin{enumerate}
\item Any minimal prime of $S$ is a homogeneous ideal of $S$.
\item Given a homogeneous ideal $I \subset S$ any minimal
prime over $I$ is homogeneous.
\end{enumerate}
\end{lemma}

\begin{proof}
The first assertion holds because the prime $\mathfrak q$ constructed in
Lemma \ref{lemma-smear-out} satisfies $\mathfrak q \subset \mathfrak p$.
The second because we may consider $S/I$ and apply the first part.
\end{proof}

\begin{lemma}
\label{lemma-dehomogenize-finite-type}
Let $R$ be a ring. Let $S$ be a graded $R$-algebra. Let $f \in S_{+}$
be homogeneous. Assume that $S$ is of finite type over $R$. Then
\begin{enumerate}
\item the ring $S_{(f)}$ is of finite type over $R$, and
\item for any finite graded $S$-module $M$ the module $M_{(f)}$
is a finite $S_{(f)}$-module.
\end{enumerate}
\end{lemma}

\begin{proof}
Choose $f_1, \ldots, f_n \in S$ which generate $S$ as an $R$-algebra.
We may assume that each $f_i$ is homogeneous (by decomposing each $f_i$
into its homogeneous components). An element of $S_{(f)}$ is a sum
of the form
$$
\sum\nolimits_{e\deg(f) =
\sum e_i\deg(f_i)} \lambda_{e_1 \ldots e_n} f_1^{e_1} \ldots f_n^{e_n}/f^e
$$
with $\lambda_{e_1 \ldots e_n} \in R$. Thus $S_{(f)}$ is generated
as an $R$-algebra by the $f_1^{e_1} \ldots f_n^{e_n} /f^e$ with the
property that $e\deg(f) = \sum e_i\deg(f_i)$. If $e_i \geq \deg(f)$
then we can write this as
$$
f_1^{e_1} \ldots f_n^{e_n}/f^e =
f_i^{\deg(f)}/f^{\deg(f_i)} \cdot
f_1^{e_1} \ldots f_i^{e_i - \deg(f)} \ldots f_n^{e_n}/f^{e - \deg(f_i)}
$$
Thus we only need the elements $f_i^{\deg(f)}/f^{\deg(f_i)}$ as well
as the elements $f_1^{e_1} \ldots f_n^{e_n} /f^e$ with
$e \deg(f) = \sum e_i \deg(f_i)$ and $e_i < \deg(f)$.
This is a finite list and we see that (1) is true.

\medskip\noindent
To see (2) suppose that $M$ is generated by homogeneous elements
$x_1, \ldots, x_m$. Then arguing
as above we find that $M_{(f)}$ is generated as an $S_{(f)}$-module
by the finite list of elements of the form
$f_1^{e_1} \ldots f_n^{e_n} x_j /f^e$
with $e \deg(f) = \sum e_i \deg(f_i) + \deg(x_j)$ and
$e_i < \deg(f)$.
\end{proof}

\begin{lemma}
\label{lemma-homogenize}
Let $R$ be a ring.
Let $R'$ be a finite type $R$-algebra, and let $M$ be a finite $R'$-module.
There exists a graded $R$-algebra $S$, a graded $S$-module $N$ and
an element $f \in S$ homogeneous of degree $1$ such that
\begin{enumerate}
\item $R' \cong S_{(f)}$ and $M \cong N_{(f)}$ (as modules),
\item $S_0 = R$ and $S$ is generated by finitely many elements
of degree $1$ over $R$, and
\item $N$ is a finite $S$-module.
\end{enumerate}
\end{lemma}

\begin{proof}
We may write $R' = R[x_1, \ldots, x_n]/I$ for some ideal $I$.
For an element $g \in R[x_1, \ldots, x_n]$ denote
$\tilde g \in R[X_0, \ldots, X_n]$ the element homogeneous of minimal
degree such that $g = \tilde g(1, x_1, \ldots, x_n)$.
Let $\tilde I \subset R[X_0, \ldots, X_n]$ generated by all
elements $\tilde g$, $g \in I$.
Set $S = R[X_0, \ldots, X_n]/\tilde I$ and denote $f$ the image
of $X_0$ in $S$. By construction we have an isomorphism
$$
S_{(f)} \longrightarrow R', \quad
X_i/X_0 \longmapsto x_i.
$$
To do the same thing with the module $M$ we choose a presentation
$$
M = (R')^{\oplus r}/\sum\nolimits_{j \in J} R'k_j
$$
with $k_j = (k_{1j}, \ldots, k_{rj})$. Let $d_{ij} = \deg(\tilde k_{ij})$.
Set $d_j = \max\{d_{ij}\}$. Set $K_{ij} = X_0^{d_j - d_{ij}}\tilde k_{ij}$
which is homogeneous of degree $d_j$. With this notation we set
$$
N = \Coker\Big(
\bigoplus\nolimits_{j \in J} S(-d_j) \xrightarrow{(K_{ij})} S^{\oplus r}
\Big)
$$
which works. Some details omitted.
\end{proof}




